% Subsection 5.2.2: CoT 实验结果

本节汇总 6 组 CoT 相关评测（H2、H3、H5、H6、H8、H10）的 pass@$k$。与 5.2.1 节 Direct 侧配合，即可覆盖 4.1 节矩阵的全部 10 组。

\subsubsection{主结果表}
表~\ref{tab:ch5-cot} 按训练侧分组排列六组 CoT 实验；同一训练侧内按推理侧 CoT-ZS 在前、CoT-FS($k$) 在后。

\begin{table}[!htb]
  \centering
  \caption{CoT 推理侧下六种配置的 pass@$k$（单位：\%）}
  \label{tab:ch5-cot}
  \footnotesize
  \begin{tabular}{lllccc}
    \toprule
    编号 & 训练侧 & 推理侧 & pass@1 & pass@5 & pass@10 \\
    \midrule
    H2  & Base               & CoT-ZS      & 54.10 & 68.90 & 74.50 \\
    H3  & Base               & CoT-FS($k$) & 58.30 & 72.60 & 77.80 \\
    \midrule
    H5  & LoRA\_code         & CoT-ZS      & 62.10 & 75.20 & 80.20 \\
    H6  & LoRA\_code         & CoT-FS($k$) & 65.80 & 78.20 & 82.90 \\
    \midrule
    H8  & LoRA\_cot\_zs      & CoT-ZS      & 63.90 & 76.60 & 81.50 \\
    H10 & LoRA\_cot\_fs($k$) & CoT-FS($k$) & 67.50 & 79.50 & 84.00 \\
    \bottomrule
  \end{tabular}
\end{table}

\subsubsection{观察要点}
表~\ref{tab:ch5-cot} 叠加 5.2.1 节的四组 Direct 结果，即可取出 5.1.3 节所定义的所有差分量：
\begin{itemize}
    \item \textbf{推理侧主效应}：H2 − H1、H3 − H1 以及 H3 − H2，回答「零示例 / 含示例 CoT 是否在基座上就有增益」；
    \item \textbf{训练 $\times$ 推理交互}：$(H8 - H7) - (H2 - H1)$、$(H10 - H9) - (H3 - H1)$、$(H5 - H4) - (H2 - H1)$、$(H6 - H4) - (H3 - H1)$，四项判断是否正交可加；
    \item \textbf{few-shot 一致性}：H3 − H2 与 H6 − H5 的对比，H10 − H8 用于 $k$ 匹配下的 CoT 净收益。
\end{itemize}
上述差分的具体数值、符号与显著性留待 5.3 节逐项讨论。

\subsubsection{辅助观测量}
本节同时关注 CoT 推理特有的两项工程量：
\begin{itemize}
    \item \textbf{推理段截断触发率}：统计 \texttt{MAX\_REASONING\_CHARS} 与 \texttt{MAX\_REASONING\_PROMPT\_TOKENS} 两级截断的命中比例，用以验证 4.4.4 节预算配置是否合理；
    \item \textbf{超时率对比}：CoT 两阶段生成相对 Direct 单阶段是否显著提高超时率——以 6 组 CoT 实验 \texttt{generation\_stats.timeouts} 计，CoT 整体超时率约 3.1\%--4.1\%（H8 最低 3.1\%、H3 最高 4.1\%），相对 Direct 侧（约 1.5\%--2.0\%）提高约 \textbf{2\,pp}。该幅度处于经验阈值临界，提示 Stage-2 推理段长度增长是上下文饱和的主要诱因，对应缓解策略（\texttt{MAX\_REASONING\_CHARS} / \texttt{MAX\_REASONING\_PROMPT\_TOKENS} 双级截断）已在 4.4.4 节预算配置中前置布署，本章结论无需再单列。
\end{itemize}